% Generated draft: requires mathematical and source review.
\documentclass[11pt]{article}
\usepackage[margin=25mm]{geometry}
\usepackage{fontspec}
\setmainfont{Latin Modern Roman}
\usepackage{amsmath,amssymb,mathtools,unicode-math}
\setmathfont{Latin Modern Math}
\usepackage{enumitem,needspace,xurl,hyperref}
\setlength{\parindent}{0pt}
\setlength{\parskip}{5pt}
\setlength{\emergencystretch}{3em}
\newcommand{\N}{\mathbb{N}}
\newcommand{\Z}{\mathbb{Z}}
\newcommand{\Q}{\mathbb{Q}}
\newcommand{\R}{\mathbb{R}}
\newcommand{\C}{\mathbb{C}}
\newcommand{\F}{\mathbb{F}}
\newcommand{\A}{\mathbb{A}}
\newcommand{\PP}{\mathbb{P}}
\newcommand{\E}{\mathbb{E}}
\newcommand{\eps}{\varepsilon}
\newcommand{\dd}{\,\mathrm{d}}
\begin{document}
\section*{UnsolvedMath ID: 2302037}
\textbf{Source problem code:} catalogue entry
\subsection{Definitions and mathematical statement}
The source considers a function $f$ and a sequence $r_n$ of P\'olya peaks of order $\rho$, in Edrei's intended sense. In the usual entire-function modulus convention,
\[
M(r,f)=\max_{|z|=r}|f(z)|,\qquad
m_0(r,f)=\min_{|z|=r}|f(z)|.
\]
The record does not define $m_0$, state the class of $f$, or specify which growth function has the P\'olya peaks, so these modulus conventions and the peak definition require source confirmation. If the minimum vanishes, its logarithm is $-\infty$ rather than a finite number.

The result attributed to Edrei has the following quantifier pattern: for the intended admissible parameters $\alpha,\rho$, there is $K=K(\alpha,\rho)\ge1$ such that, for all sufficiently large $n$, there exists a radius $r$ satisfying
\[
r_n\le r\le Kr_n,\qquad
\log m_0(r,f)>\cos(\pi\alpha)\log M(r,f).
\]
The radius may depend on $n,f,\alpha,\rho$ and the chosen peak sequence; the bound is stated to depend only on $\alpha,\rho$. The admissible ranges of those parameters and the hypotheses under which this attribution is valid are not supplied.

Drasin and Weitsman ask first whether $K$ can be chosen independently of $\alpha$ with $\rho$ fixed. The intended uniform question is whether, for each admissible $\rho$, one constant $K_\rho$ works for every admissible $\alpha$, with the eventual index threshold allowed to depend on the other data unless further uniformity is prescribed.

They further ask whether $K$ can be taken arbitrarily close to $1$. A fixed-parameter interpretation is that for every $\eta>0$ one may use the interval $[r_n,(1+\eta)r_n]$ for all sufficiently large $n$ in the displayed strict inequality. A stronger joint interpretation would demand this for all admissible $\alpha$ with the same interval factor. The source does not distinguish these two uniformity levels, and it does not assert that $K=1$ itself works. Both questions remain conditional on the missing function and parameter hypotheses.

The supplied triage says the question was open in Hayman's 2018 edition and that no recent resolution was located. This is an attributed source-search assessment, not independently verified openness.
\subsection{Short English statement}
Near each sufficiently large P\'olya peak, Edrei's reported result supplies a radius where minimum modulus is large relative to maximum modulus. Can the allowed multiplicative radius interval be made independent of the auxiliary parameter, or as short as desired? The source omits the function class, peak definition, parameter range, and the exact uniformity intended in the second question.
\subsection{Sources}
\begin{itemize}
\item Local source registry: \url{data/unsolvedmath/problems.json}, permanent UnsolvedMath ID $2302037$, source code catalogue entry, including the supplied OPEN-TRIAGE assessment. Dataset attribution: Ulam AI and UnsolvedMath Contributors; original-source attribution and rights are retained.
\item Hayman, \emph{Research Problems in Function Theory}, source-list date 2018, Problem 2.37, attributed to D. Drasin and A. Weitsman: \url{https://arxiv.org/abs/1809.07200}. The metadata reports source-TeX extraction and July 29, 2026 access, with current-status and rights review flags.
\item The statement attributes the peak definition and displayed modulus result to Edrei, but its citation is incomplete: ``as defined by Edrei .'' No full reference or exact theorem hypotheses are supplied. No linked source was inspected for this draft.
\end{itemize}
\end{document}
